\clearpage\pagestyle{otMain}\setlength{\otBlocEspacement}{2.5mm}\fontsize{11}{15}\selectfont\corrigesChapitre{14}{Durées : jusqu’à la seconde}
\begin{blocLivre}{cor-14-1}\begin{corrige}{1}{Convertir en secondes}\demarche{Compter 60 secondes par minute.}
$2\times60=\res{120}$ s ; $120+15=\res{135}$ s ; $5\times60+8=\res{308}$ s.
\end{corrige}\end{blocLivre}
\begin{blocLivre}{cor-14-2}\begin{corrige}{2}{Convertir en minutes et secondes}\demarche{Extraire les paquets de 60.}
$200=3\times60+20$ : $\res{3\ \text{min}\ 20\ \text s}$. $428=7\times60+8$ : $\res{7\ \text{min}\ 08\ \text s}$. 59 s reste 59 s ; 60 s font $\res{1\ \text{min}}$ exactement.
\end{corrige}\end{blocLivre}
\begin{blocLivre}{cor-14-3}\begin{corrige}{3}{L’heure en secondes}\demarche{Multiplier les paquets de 60.}
$60\times60=\res{3\,600}$ s. Un quart d’heure : $3\,600:4=\res{900}$ s, soit 15 minutes.
\end{corrige}\end{blocLivre}
\begin{blocLivre}{cor-14-4}\begin{corrige}{4}{Additionner des durées}\demarche{Échanger 60 unités contre une du rang supérieur.}
$2$ min 40 s $+1$ min 35 s $=3$ min 75 s $=\res{4\ \text{min}\ 15\ \text s}$. Et $3$ h 45 $+2$ h 30 $=5$ h 75 min $=\res{6\ \text h\ 15}$.
\end{corrige}\end{blocLivre}
\begin{blocLivre}{cor-14-5}\begin{corrige}{5}{Le chrono du 100 m}\demarche{Tout convertir en secondes.}
$1$ min 01 s $=61$ s : $\res{\text{Noa}}$ gagne, avec $61-58=\res3$ s d’avance.
\end{corrige}\end{blocLivre}
\begin{blocLivre}{cor-14-6}\begin{corrige}{6}{La durée du trajet}\demarche{Faire étape à l’heure ronde.}
Jusqu’à 11 h : 12 min 30 s ; puis 2 min 10 s : total $\res{14\ \text{min}\ 40\ \text s}$.
\end{corrige}\end{blocLivre}
\begin{blocLivre}{cor-14-7}\begin{corrige}{7}{L’heure de fin}\demarche{Compléter d’abord la minute ou l’heure entamée.}
$9$ h $52+8$ min $=10$ h, puis 17 min : $\res{10\ \text h\ 17}$. Pour la machine : 30 s mènent à 17 h 49, puis 1 min 10 s : $\res{17\ \text h\ 50\ \text{min}\ 10\ \text s}$.
\end{corrige}\end{blocLivre}
\begin{blocLivre}{cor-14-8}\begin{corrige}{8}{Remonter le temps}\demarche{Soustraire par étapes depuis la fin.}
On recule : 10 min ramènent à 18 h, puis 25 min : enfourner à $\res{17\ \text h\ 35}$. Contrôle : $17$ h $35+35$ min $=18$ h 10.
\end{corrige}\end{blocLivre}
\begin{blocLivre}{cor-14-9}\begin{corrige}{9}{Comparer des durées}\demarche{Convertir dans la même unité.}
$1$ min 30 s $=90$ s $<95$ s : $\res{95\ \text s}$ gagne. $2$ h 05 $=125$ min $>118$ min : $\res{2\ \text h\ 05}$ est plus longue.
\end{corrige}\end{blocLivre}
\begin{blocLivre}{cor-14-10}\begin{corrige}{10}{Les dixièmes qui décident}\demarche{Comparer les décimales des secondes.}
$32{,}48<32{,}60$ : $\res{\text{Lou}}$ gagne. Écart : $32{,}6-32{,}48=\res{0{,}12}$ s, douze centièmes de seconde.
\end{corrige}\end{blocLivre}
\begin{blocLivre}{cor-14-11}\begin{corrige}{11}{La playlist}\demarche{Multiplier en secondes puis reconvertir.}
$3$ min 45 s $=225$ s ; $225\times4=900$ s $=\res{15\ \text{min}}$ exactement : le quart d’heure est atteint, $\res{\text{pas dépassé}}$.
\end{corrige}\end{blocLivre}
\begin{blocLivre}{cor-14-12}\begin{corrige}{12}{Le tour de circuit}\demarche{Travailler en secondes.}
$72\times5=360$ s $=\res{6\ \text{min}}$. En 600 s : $600=72\times8+24$, soit $\res8$ tours complets et 24 s entamées.
\end{corrige}\end{blocLivre}
\begin{blocLivre}{cor-14-13}\begin{corrige}{13}{La séance bien remplie}\demarche{Soustraire les parties connues du total.}
$1$ h 30 $=90$ min ; $90-12-55=\res{23}$ min d’étirements.
\end{corrige}\end{blocLivre}
\begin{blocLivre}{cor-14-14}\begin{corrige}{14}{Le journal de bord}\demarche{Calculer chaque tronçon puis additionner.}
Matin : de 9 h 40 à 12 h 05, soit 2 h 25. Après-midi : de 12 h 50 à 15 h 10, soit 2 h 20. Total : $2$ h $25+2$ h $20=\res{4\ \text h\ 45}$ de marche.
\end{corrige}\end{blocLivre}
\begin{blocLivre}{cor-14-15}\begin{corrige}{15}{Vrai ou faux ?}\demarche{Échanger les paquets de 60, pas de 100.}
$85$ min dépassent l’heure : 1 h 85 min $=\res{2\ \text h\ 25}$. Et 30 s font une demi-minute : 2 min 30 s $=\res{2{,}5}$ min, pas $2{,}3$. Le temps compte en base 60.
\end{corrige}\end{blocLivre}
